\documentclass[11pt]{article}
\usepackage[T1]{fontenc}
\usepackage[utf8]{inputenc}
\usepackage[a4paper,margin=0.92in]{geometry}
\usepackage{amsmath,amssymb,amsthm,mathtools,booktabs,graphicx,natbib,enumitem}
\usepackage[hidelinks]{hyperref}
\usepackage{microtype}
\usepackage{fancyhdr}
\pagestyle{fancy}\fancyhf{}\fancyhead[L]{\small Portfolio-aware alpha retirement and revival}\fancyhead[R]{\small Working paper}\fancyfoot[C]{\thepage}
\setlength{\emergencystretch}{1.5em}
\setlength{\headheight}{14pt}\setlength{\parskip}{3pt}
\newtheorem{assumption}{Assumption}\newtheorem{theorem}{Theorem}\newtheorem{proposition}{Proposition}\newtheorem{corollary}{Corollary}
\newcommand{\E}{\mathbb E}\newcommand{\Prob}{\mathbb P}\newcommand{\F}{\mathcal F}\newcommand{\ind}{\mathbf 1}\newcommand{\clip}{\operatorname{clip}}\newcommand{\Cov}{\operatorname{Cov}}\newcommand{\Var}{\operatorname{Var}}\newcommand{\KL}{\operatorname{KL}}\newcommand{\kl}{\operatorname{kl}}
\title{\vspace{-1.2cm}When Should I Kill---and Revive---an Alpha?\\[0.3em]\large Portfolio-Aware Sequential Evidence and the Cost of Certification}
\author{Enrico Bignozzi}
\date{6 October 2026\\\small Research draft with executed simulations and JKP historical replay}
\begin{document}
\maketitle\thispagestyle{empty}
\begin{abstract}
A strategy can have a negative expected return and still improve a portfolio through diversification. Conversely, a profitable strategy can reduce the utility of a particular allocation. We study reversible decisions to park and reactivate fixed strategy allocations using paired, predictable portfolio comparisons. Bounded utility-difference scores feed directional betting processes, restart mixtures, and an explicit error budget across strategies and decision episodes. The resulting control concerns false alarms while an episode's conditional-mean null remains true; it does not certify the instantaneous sign of future economic value. We establish this guarantee, a persistent-signal delay bound, a canonical information lower bound, and the relation between clipped and raw utility targets. All proofs appear in the appendix. Executed simulations distinguish economic gains from statistical conservatism. The method preserves negative-return hedges and benefits from revival after strong mean deterioration, but it reacts too slowly in a correlation-only revival design and underperforms an always-on policy there. A JKP-only replay of 153 US factors from January 2000 through December 2025 produces no certified switches under the reference specification. These negative results identify the main obstacle to deployment: demanding lifetime error control can exhaust the economic opportunity before sufficient evidence accumulates. The package provides an auditable research baseline, not evidence of a state-of-the-art profitable allocator.
\end{abstract}
\noindent\textbf{Keywords:} portfolio choice; sequential testing; e-processes; strategy lifecycle; shadow portfolios; nonstationarity.\par
\noindent\textbf{JEL:} C12, C14, C58, G11.
\vspace{0.5em}
\noindent\textbf{Scope.} No orders are submitted. JKP costs are transparent sensitivity assumptions, not estimated stock-level execution costs. The historical replay respects observation timing but does not reconstruct historical publication dates or data vintages.
\newpage

\section{Introduction}
The decision to stop funding an investment strategy is distinct from the decision to reject a zero-alpha hypothesis. Capital is allocated to a book, strategies interact through covariance, and a temporarily unattractive strategy may recover. This paper treats retirement as a reversible allocation decision rather than an irreversible declaration that an economic effect has disappeared.

The central object is the difference in realized utility between two portfolios that could both have been specified before observing the next return: one maintains a candidate allocation and the other replaces that allocation with cash. Their difference automatically includes the candidate's interaction with the rest of the book. When the candidate is parked, its return remains observable in a shadow portfolio. The same comparison therefore supports subsequent reactivation.

This construction requires two qualifications. First, an optional asset cannot reduce the population optimum when its weight may be set to zero. A negative contribution must refer to an implemented allocation, a constrained policy, or a genuine cost of maintaining access. We use fixed capital slots, so that parking a strategy changes an actual feasible decision. Second, continuous monitoring does not make inference about arbitrarily changing future conditional means possible. Our statistical guarantee is an episode-wise false-alarm guarantee, not a certificate that every switch is economically correct at the instant it occurs.

Our methodological contribution is a complete, auditable specification of this decision problem. The elementary betting factors, mixtures, and maximal-inequality arguments are not claimed as new statistical devices. We combine them with paired portfolio comparisons, reversible monitoring, explicit spending across strategies and episodes, and safe finite-memory implementation. This yields a transparent baseline on which stronger economic and statistical designs can be assessed.

The executed evidence is mixed. A negative-return diversifier remains active under the certified rule, while standalone rolling criteria repeatedly remove it. Strong mean deterioration is detected and revival improves utility relative to permanent retirement. However, the key correlation-only revival experiment exposes a large delay cost. The reference certified rule loses 3.38 units of $10^{-4}$ utility per period relative to always on, whereas the portfolio rolling comparator gains 20.65. In the historical JKP replay the certified policies never switch. We report these outcomes rather than treating nonrejection as successful active management.

\subsection{Relation to existing work}
Time-uniform confidence sequences formalize inference at data-dependent stopping times \citep{howard2021}. Betting constructions provide practical inference for bounded observations \citep{waudby2024}. Our score transformation deliberately enters this bounded setting. This route does not require exchangeable conformal ranks, but it does require a conditional moment restriction under the null.

Restarted evidence also appears in sequential change detection. \citet{shin2022} construct e-detectors and control average run length. A unit-capital mixture controlling the probability of ever crossing a boundary is a different error criterion; the two guarantees must not be interchanged. \citet{bhattacharyya2026} study optimal conformal change detection for independent observations with unknown pre- and post-change laws. We do not transfer their minimax conclusion to financial returns or to our reversible policy.

Return-to-baseline monitoring \citep{regis2026} and sequential tradeability testing \citep{stephan2026} are closely related to recovery and deployment decisions. More importantly, \citet{bonacorsi2026} explicitly studies how certification time consumes a decaying alpha's remaining economic capacity. Therefore, neither statistical certification of alpha nor the economic cost of waiting is a first contribution here. The specific object is the reversible value of an allocation relative to a changing portfolio, with a complete paired-book experiment.

Our empirical source is the factor library of \citet{jkp2023}. The current official portal supplies signed factor returns and theme aggregates \citep{jkpdata}. These are suitable for a reproducible historical replay, but a contemporary retrospective library cannot by itself establish historical investability of every included signal.

\section{The economic object: a policy, not an opportunity set}\label{sec:economy}
Let $\F_{t-1}$ contain all observations and decisions available before period $t$. There is a core return $R_{0,t}$ and $M$ candidate strategy returns $R_{1,t},\ldots,R_{M,t}$. All returns are measured relative to cash. A state $a_{j,t}\in\{0,1\}$ is chosen using $\F_{t-1}$, with one denoting an active allocation. Candidate $j$ receives a fixed slot $q_j>0$. The reference experiment uses $w_0=0.8$ and $q_j=0.2/M$. Parked capital remains in cash; other strategies are not scaled up.

With a predictable recurring deduction $f_{j,t}$, define
\begin{equation}
 A_{j,t}=q_j(R_{j,t}-f_{j,t}),\qquad
 B_{j,t}=w_0R_{0,t}+\sum_{i\ne j}a_{i,t}A_{i,t}.
\end{equation}
The paired books are $R^{(+j)}_t=B_{j,t}+A_{j,t}$ and $R^{(-j)}_t=B_{j,t}$. They share all other positions and use the same realized return vector. The incremental realized utility is
\begin{equation}\label{eq:paired}
 D_{j,t}=u_\gamma(B_{j,t}+A_{j,t})-u_\gamma(B_{j,t}),
 \qquad u_\gamma(r)=r-\frac{\gamma}{2}r^2.
\end{equation}
This is expected quadratic utility, not exactly the mean--variance certainty equivalent. Specifically, $\E u_\gamma(R)=\E R-\gamma\Var(R)/2-\gamma(\E R)^2/2$. We report certainty equivalents separately where appropriate.

\begin{proposition}[Exact marginal utility]\label{prop:marginal}
Let conditional second moments exist. Write $m_A=\E[A_{j,t}\mid\F_{t-1}]$, $m_B=\E[B_{j,t}\mid\F_{t-1}]$, and $v_A=\Var(A_{j,t}\mid\F_{t-1})$. Then
\begin{equation}\label{eq:conditionalvalue}
 \E[D_{j,t}\mid\F_{t-1}]
 =m_A-\gamma\left\{m_Am_B+\Cov(A_{j,t},B_{j,t}\mid\F_{t-1})
       +\tfrac12(v_A+m_A^2)\right\}.
\end{equation}
Consequently, $m_A<0$ is compatible with positive incremental utility.
\end{proposition}

The covariance term is the reason a standalone return test is insufficient. A negative covariance can compensate for a modestly negative mean. A positively correlated allocation can instead have negative incremental utility despite a positive mean. Neither statement implies that the asset itself should be excluded from an unconstrained investment opportunity set.

\begin{proposition}[Optional-asset monotonicity]\label{prop:optional}
If feasible portfolios without strategy $j$ are contained in the feasible set with access to $j$, and access has no unavoidable cost, then the supremum of expected utility with access is no smaller than the supremum without access.
\end{proposition}

For this reason, our simulation label ``positive redundant'' refers to a positive-return \emph{fixed allocation} that is costly in utility terms. It does not claim negative value for the option of trading that strategy. A leave-one-out optimizer that allows zero weights would require a different estimand, such as the value of a fitted finite-sample policy including estimation and maintenance costs.

\subsection{Trading costs and observation}
The actually funded book deducts a switching charge
\begin{equation}
 C_t^{\rm sw}=c_{\rm sw}\sum_jq_j|a_{j,t}-a_{j,t-1}|.
\end{equation}
The paired scores include recurring costs but not the full future cost of switching and possible switching back. Hysteresis is a friction-control heuristic, not a solved optimal stopping problem. A cash opportunity cost can be added to $f_{j,t}$ when returns are not already relative to the relevant alternative.

Shadow returns must remain observable without placing real orders. Factor return panels meet this accounting requirement, but not the stronger requirement that hypothetical fills, borrow availability, or price impact equal real execution. Our empirical analysis does not make that stronger identification claim.


\section{Bounded evidence and reversible decisions}\label{sec:method}
Choose a strictly positive, predictable scale $s_{j,t}$ and define
\begin{equation}\label{eq:score}
 Z_{j,t}=\clip\left(D_{j,t}/s_{j,t},-1,1\right),\qquad
 m_{j,t}=\E[Z_{j,t}\mid\F_{t-1}].
\end{equation}
The reference scale is fixed after 60 burn-in observations:
$s_{j,t}=q_j\max(\widehat\sigma_{j,\rm burn},0.005)$.
Thus future returns cannot enter normalization. Clipping controls tail influence but changes the statistical target. A scale estimated from the same period's return would violate the stated predictability requirement.

Let $h\in[0,1)$ be a symmetric indifference band. During an active episode the null is $m_{j,t}\ge-h$; during a parked episode it is $m_{j,t}\le h$. The respective evidence inputs are
\begin{equation}\label{eq:input}
 X^-_{j,t}=\frac{-Z_{j,t}-h}{1+h},\qquad
 X^+_{j,t}=\frac{Z_{j,t}-h}{1+h}.
\end{equation}
They lie in $[-1,1]$ and have nonpositive conditional mean under the corresponding null. The band avoids demanding a reversal whenever a noisy score crosses zero. It does not imply a fixed dollar cost threshold because the scale maps utility into normalized units.

\subsection{Directional betting and restart mixture}
For $0<\lambda<1$, $1+\lambda X_{j,t}$ is nonnegative and has conditional expectation at most one under the episode null. A process launched at time $s_\ell$ is
\begin{equation}
 K_{\ell,\lambda,t}=\prod_{v=s_\ell}^{t}(1+\lambda X_{j,v}).
\end{equation}
All starts are predictable. We launch every $d=24$ observations within an episode and mix seven fractions
$\Lambda=\{1/32,1/16,1/8,1/4,1/2,3/4,0.9\}$ with masses $p_\lambda=1/7$.
Start $\ell$ has mass $\pi_\ell=1/[\ell(\ell+1)]$. If $L_t$ starts have occurred, the evidence is
\begin{equation}\label{eq:mixture}
 E_{j,t}=\underbrace{\frac{1}{L_t+1}}_{\text{unstarted capital}}
 +\sum_{\substack{\ell\le L_t:\ t-s_\ell<H}}
       \pi_\ell\sum_{\lambda\in\Lambda}p_\lambda K_{\ell,\lambda,t},
 \qquad H=720.
\end{equation}
Expired components are discarded, not redistributed. The unstarted capital is essential: omitting it and interpreting a truncated sum as a unit-initial-capital martingale changes the accounting. Discarding old wealth reduces evidence and is safe for false-alarm control, although it can be costly for power.

\subsection{Error spending across strategies and episodes}
Number each strategy's episodes $k=1,2,\ldots$, including both active and parked states. Allocate
\begin{equation}\label{eq:spending}
 \alpha_{j,k}=\frac{\delta}{M k(k+1)},\qquad
 \sum_{j=1}^{M}\sum_{k=1}^{\infty}\alpha_{j,k}=\delta.
\end{equation}
The decision threshold is $1/\alpha_{j,k}$. A fresh state begins with a fresh unit-capital process and a \emph{smaller} error budget. Resetting to a 5\% test after every switch would not give a 5\% lifetime family-wise guarantee.

At most one candidate is switched per book per period: among candidates crossing their boundary, accept the largest evidence-to-threshold ratio. Implement that switch at $t+1$, not at $t$. Its new direction then begins a new episode. Other tests continue with the resulting predictable book. The latter convention means their target is an adaptive reference-book comparison, not a fixed intrinsic property of a factor. Atomic switching avoids applying mutually inconsistent simultaneous leave-one-out recommendations; it is not an optimization theorem.

\subsection{Operational algorithm}
\begin{enumerate}[leftmargin=1.5em,itemsep=1pt]
\item Before observing returns, store the active set, paired allocation rules, score scales, and applicable error budgets.
\item Observe all strategy returns, including shadow returns for parked strategies. Compute the paired utilities in \eqref{eq:paired} and the bounded scores in \eqref{eq:score}.
\item Update directional evidence using \eqref{eq:input}--\eqref{eq:mixture}. Log recurring deductions, actual switching charges, and unclipped utility separately.
\item Accept at most one threshold-crossing decision, change its state for the next period, and increment that strategy's episode counter. Otherwise keep the current book.
\end{enumerate}
No fitted covariance matrix, return-forecasting network, or latent-state estimator is required by this reference implementation. Such models may specify predictable comparators or betting fractions, but they cannot turn a violated conditional null into a valid one.

\section{What is guaranteed---and what is not}\label{sec:theory}
\begin{assumption}[Predictability and bounded evidence]\label{ass:bounded}
For each episode, paired portfolio rules, scales, starts, and betting fractions are predictable. Observed evidence inputs lie in $[-1,1]$. Before the first violation of that episode's null, $\E[X_{j,t}\mid\F_{t-1}]\le0$.
\end{assumption}
This formulation permits heteroskedasticity, serial dependence, and cross-strategy dependence. It does not say that arbitrary dependent financial returns satisfy the null. The restriction is conditional on the \emph{full} information set, not merely that the unconditional mean is nonpositive. No independence or mixing claim is needed for validity once the stated conditional restriction holds.

\begin{theorem}[Lifetime false-alarm control on valid prefixes]\label{thm:valid}
Under Assumption~\ref{ass:bounded}, let $\sigma_{j,k}$ be the start of episode $k$ for strategy $j$, and let $\nu_{j,k}$ be the first time in that episode at which its directional conditional-mean null fails. Set $\nu_{j,k}=\infty$ if it never fails. With \eqref{eq:mixture}--\eqref{eq:spending},
\begin{equation}\label{eq:fwer}
 \Prob\left(\text{at least one accepted switch in an episode }(j,k)
                  \text{ at a time }t<\nu_{j,k}\right)\le\delta.
\end{equation}
The result allows data-dependent episode starts and arbitrary dependence across strategies. Discarding expired wealth or limiting accepted switches can only reduce this crossing event.
\end{theorem}
An episode that previously experienced a null violation can contain stale evidence. If its economic environment reverses before a delayed alarm, \eqref{eq:fwer} does not count that alarm as a false alarm on an uninterrupted valid prefix. In particular, \emph{it is incorrect to restate this theorem as a 95\% guarantee that every retirement or revival matches the current sign of economic value}. The distinction is material in the correlation-revival experiment.

\begin{theorem}[Persistent-signal sufficient delay]\label{thm:delay}
Consider one episode and an available restart time $s_\ell$. Suppose that until an accepted switch or for the next $n$ observations, whichever is earlier,
$\E[X_{j,t}\mid\F_{t-1}]\ge\Delta>0$.
Assume there is a grid fraction $\lambda\in[\Delta/4,\Delta/2]$, the component is not expired during these observations, and no numerical saturation prevents its boundary crossing. Write
\begin{equation}
 A=\log\frac{1}{\alpha_{j,k}\pi_\ell p_\lambda},\qquad
 n\ge\frac{128}{\Delta^2}\left(A+\log\frac1\beta\right),\qquad 0<\beta<1.
\end{equation}
Then its evidence crosses $1/\alpha_{j,k}$ by the end of these $n$ observations with probability at least $1-\beta$, unless that strategy has already switched. For an isolated strategy this implies an accepted switch. With competing threshold crossings, acceptance additionally requires that the atomic decision rule does not defer the strategy.
\end{theorem}
This is a conservative sufficient bound, not an asymptotic equality, an unconditional expected-delay bound, or a claim of minimax optimality. It requires a sufficiently long regime and $n\le H$. An unknown onset can add up to $d-1$ observations to reach the next scheduled start. Because $A$ includes the restart and episode budgets, late changes and repeated decisions become more expensive to certify. A fixed finite grid also limits small-signal adaptation; the stated interval need not contain a grid point for arbitrarily small $\Delta$.

\begin{proposition}[A canonical information obstruction]\label{prop:lower}
For independent observations in $\{-1,1\}$, let $P_0(X=1)=1/2$ and $P_\Delta(X=1)=(1+\Delta)/2$, where $0<\Delta\le1/2$. If a decision based on at most $n$ observations has null probability at most $\alpha$ and power at least $1-\beta>\alpha$, then
\begin{equation}
 n\ge\frac{\kl(1-\beta,\alpha)}{\KL(P_\Delta\Vert P_0)}
 \ge\frac{\kl(1-\beta,\alpha)}{\Delta^2}.
\end{equation}
\end{proposition}
This lower bound concerns a canonical fixed-horizon subproblem. It does not prove a matching lower bound for the unknown-change, many-strategy, reversible policy with transaction costs.

\begin{proposition}[Clipping changes the target]\label{prop:clip}
For a predictable $s>0$, $Z=\clip(D/s,-1,1)$, and conditional integrability,
\begin{equation}
 |\E[D\mid\F]-s\E[Z\mid\F]|
 \le\E[(|D|-s)_+\mid\F].
\end{equation}
If additionally $\E[|D|^{1+\eta}\mid\F]\le V$ for some $\eta>0$, then the right side is at most $V/s^\eta$.
\end{proposition}
A raw-utility sign statement requires control of this remainder. Our empirical implementation does not assume a known financial tail-moment bound and therefore makes no such conversion.

\begin{proposition}[No unrestricted instantaneous future classifier]\label{prop:impossibility}
Suppose two admissible data laws agree on $\F_t$ but assign opposite signs to the next-period conditional incremental utility. For any binary $\F_t$-measurable allocation decision, the larger of its two misclassification probabilities is at least $1/2$.
\end{proposition}
Persistence, predictability restrictions on future regimes, abstention, or an explicit decision-theoretic loss are necessary to strengthen the interpretation. The paper's guarantees do not solve that separate problem.


\section{Executed simulation study}\label{sec:sim}
\subsection{Design and estimands}
All policies see the same return paths within each simulation setting. The common-factor construction is
\begin{align}
 R_{0,t}&=0.006+s_tF_t,\\
 R_{j,t}&=\mu_{j,t}+s_t\{\rho_{j,t}F_t+(1-\rho_{j,t}^2)^{1/2}\varepsilon_{j,t}\}.
\end{align}
For the Gaussian reference design, $s_t=0.08$ and all innovations are independent standard normals. This construction gives a positive-semidefinite covariance matrix for every regime. Conditional means and covariances are known to the evaluator but not to the competing algorithms.

Seven designs cover permanent mean death; mean death and revival; a negative-return hedge; a positive-return redundant allocation; a constant-mean correlation reversal; gradual correlation changes; and repeated correlation cycles. Exact parameters are in Appendix~\ref{app:dgp}. The main design has 1,800 observations, 60 burn-in observations, one candidate strategy, 200 paths per scenario, and both Gaussian and standardized Student-$t_5$ innovations. These periods are simulation units, not a claim of 150 years of realistic monthly returns. In particular, the $\pm4.5\%$ mean-change designs are strong-signal diagnostics, not JKP calibrations.

Separate stress tests use GARCH volatility and an autoregressive common factor. Scaling experiments use $M\in\{20,50,100,150\}$ with heterogeneous regime phases. An additional weak-correlation design reduces correlation magnitudes by a factor of four. A dedicated delay study uses 500 paths per setting. Null calibration is conducted separately on bounded scores with known conditional symmetry, rather than assuming that economic labels imply the clipped-score null.

The reference parameters are $\gamma=8$, $\delta=0.05$, $h=0.02$, a recurring deduction of 5 basis points per unit allocation per period, and a switching charge of 10 basis points per unit allocation changed. These are pre-specified demonstration costs, not empirical estimates.

\subsection{Comparators and fairness}
Eight policies are evaluated: always on; a 60-period standalone-score rule; a 60-period paired-portfolio-score rule; a fixed CUSUM heuristic; a fixed two-state HMM filter; a cumulative e-process without restarts; restart-based permanent retirement; and restart-based retirement plus revival. All share allocation accounting and costs. The CUSUM and HMM are transparent heuristics, not fitted oracles. Their thresholds and the rolling rules are \emph{not matched to the certified policy's lifetime false-alarm rate}. Consequently, faster adaptation by a heuristic is economically relevant but not a like-for-like statistical efficiency comparison.

We do not present the current implementation as an empirical state-of-the-art comparison against optimized conformal detectors, online multiple-testing procedures, or optimal switching policies. Such a claim would require faithful implementations, matched error objectives, and tuning on disjoint simulations. The included benchmarks establish the mechanism and diagnose the reference policy, not a universal ranking.

\subsection{Performance measures}
The primary economic statistic is the path-level paired difference
\begin{equation}
 \widehat G=\frac{1}{T-T_0}\sum_{t>T_0}
       \{u_\gamma(R_t^{\rm policy})-u_\gamma(R_t^{\rm always})\}.
\end{equation}
Monte Carlo standard errors are calculated from independent path replications within a setting, not from individual serially dependent returns. The package also records actual returns, turnover through switch counts, mean activity, drawdown diagnostics, and conditional expected utility.

For $M=1$, a one-period conditional oracle chooses the better current state knowing true conditional moments, with the same previous state and switching cost. Its nonnegative gap to the method is a \emph{local opportunity-loss diagnostic}. It is not the regret to a globally optimal clairvoyant switching policy. No combinatorial oracle is fabricated for $M>1$.

Delay tables report the first correct state after a regime onset. A path already in the correct state has delay zero; undetected paths are right-censored at the regime end and assigned that cap. Detection fractions are reported beside capped means. These are economic state-adaptation diagnostics based on raw conditional utility, not empirical estimates of the theorem's bounded-score false-alarm event.

\subsection{False-alarm calibration}
\begin{table}[htbp]\centering\small
\caption{Executed null calibration: any switch over 720 periods and five strategies. The global budget is 5\%. Intervals are exact binomial 95\% intervals across 2,000 paths per null.}\label{tab:null}
\input{tables/null.tex}
\end{table}
The certified procedure switches on 7 of 2,000 bounded-independent paths and 4 of 2,000 predictable-volatility paths. Thus its observed rates, 0.35\% and 0.20\%, are well below the nominal upper bound. This supports implementation-level validity in these experiments and also reveals conservatism. It does not prove exact size or validate unrestricted real-return assumptions. A repeated unadjusted two-sided $z$ statistic produces approximately 91\% paths with at least one crossing when repeatedly examined across five strategies. That statistic is an intentionally invalid negative control, not a fair competing detector.

\begin{figure}[htbp]\centering\includegraphics[width=.77\linewidth]{../figures/null_control.pdf}
\caption{Lifetime false-alarm calibration. Error bars are exact binomial intervals; the dashed line is the upper bound, not a target frequency of profitable decisions.}\label{fig:null}
\end{figure}

\subsection{Economic results and the cost of waiting}
\begin{table}[htbp]\centering\scriptsize
\caption{Gaussian simulations: mean utility gain over always on, multiplied by $10^4$. Parentheses contain paired Monte Carlo standard errors. There are 200 paths per scenario. Positive numbers favor the policy.}\label{tab:main}
\resizebox{\linewidth}{!}{\input{tables/main_gaussian.tex}}
\end{table}

The negative-return hedge provides the clearest portfolio-awareness check. The certified method never removes it in the reference Gaussian experiment. Standalone rolling decisions instead lose 27.20 units of $10^{-4}$ utility per period. Even a portfolio rolling rule incurs a smaller loss because noisy repeated switching is not free. This is evidence that the chosen policy score captures economically relevant covariance; it is not a claim that all negative-return factors should be held.

In the strong mean-revival scenario, irreversible retirement gains 14.97 units relative to always on, whereas reversible retirement gains 27.88. The ability to redeploy therefore matters economically. The rolling portfolio comparator gains 37.32, showing that valid certification still imposes a substantial delay cost.

The correlation-only scenario is a more demanding and less favorable result. Mean returns do not change. The candidate loses and regains value through its covariance with the core. The certified reversible rule loses 3.38 utility units relative to always on, although permanent retirement loses 8.45. Revival reduces damage but does not make the method profitable relative to the passive comparator. The portfolio rolling rule gains 20.65. This result rejects a simple deployment narrative in which portfolio awareness plus an e-process is automatically sufficient.

\begin{figure}[htbp]\centering\includegraphics[width=\linewidth]{../figures/simulation_utility.pdf}
\caption{Economic comparison across Gaussian designs. Error bars show 1.96 Monte Carlo standard errors, not confidence in real-world profitability. The correlation-revival loss is retained.}\label{fig:sim}
\end{figure}

\begin{table}[htbp]\centering\small
\caption{Correlation-revival adaptation, 500 additional paths. Each regime lasts 600 observations. Capped delay includes nondetections; already-correct states have delay zero. ``Adapted'' means the desired state is reached before regime end, not that it is maintained.}\label{tab:delay}
\input{tables/delay.tex}
\end{table}
The delay study explains the utility loss. The certified rule's mean capped retirement delay is about 329 observations, versus about 30 for the portfolio rolling rule. Revival is even slower. A recovery that is shorter than the evidence-collection interval can be economically over before certification is available. Counting only successful detections would conceal this failure mode, which is why the denominator and censoring convention are explicit.

\begin{figure}[htbp]\centering\includegraphics[width=\linewidth]{../figures/correlation_path.pdf}
\caption{State paths in replication zero, selected by index rather than by appearance. Within each row, the lower level is parked and the upper level is active. Vertical markers denote covariance regime changes.}\label{fig:path}
\end{figure}

\subsection{Dependence, scale, and weak signals}
Student-$t_5$ results are retained in the full CSV output alongside Gaussian results. GARCH and autoregressive common-factor experiments use exact predictable moments for evaluation. They stress performance rather than claiming that the raw-return economic label always coincides with the certified conditional-score null.

At $M=150$, the reference correlation-revival scaling experiment has only 30 paths and is therefore a diagnostic, not high-precision evidence. The certified policy averages 14.5 switches over 840 evaluation periods and keeps about 97.5\% of slots active. Its utility difference from always on is negative. The factor-count penalty and the one-switch-per-period convention both matter; computational scalability does not imply statistical power.

In the weak-correlation design the certified rule averages 0.005 switches per path. Its practical behavior is almost identical to always on. This is the expected danger of a stringent lifetime objective when marginal signals are small. We retain weak-signal, dependence, and scaling results in the reproducibility package and report selected tables in Appendix~\ref{app:results}.


\section{JKP-only historical experiment}\label{sec:jkp}
\subsection{Source, timing, and construction}
We downloaded the official USA, monthly, capped-value-weighted JKP factor archive and its market and theme archives on 6 October 2026. The factor file contains 153 distinct names and returns through December 2025. The source manifest records original URLs, byte counts, SHA-256 hashes, and the portal's noncommercial research license. Raw files are not committed to the public repository.

January 1995--December 1999 supplies the 60-month burn-in. Factors are eligible if they have complete returns in that burn-in, without filtering on subsequent performance. All 153 are eligible, and their evaluation panel has no missing factor-months through December 2025. A future missing observation would stop the loader rather than trigger retrospective factor deletion or silent zero filling. Returns are decimal, already signed as supplied; the direction field is not multiplied a second time.

Decisions are evaluated from January 2000 through December 2025, giving 312 months. Every decision uses returns observed no later than that month and affects only the next month. Burn-in scales remain fixed. The source's current library, historical revisions, and possibly ex-post research selection remain limitations. We therefore call the exercise an \emph{out-of-sample-in-time historical replay}, not a fully point-in-time investment backtest. Factor-level data also do not identify real stock-level fills, capacity, borrow costs, or slippage.

\subsection{Reference replay}
The reference book has 80\% exposure to the JKP market and 20\% spread equally across the factor slots. This anchor makes diversification relative to a broad core visible, but market exposure dominates total volatility. Unallocated factor slots remain cash. The recurring deduction is 5 basis points per unit factor exposure per month and a change costs 10 basis points per unit exposure. Zero-cost and higher-cost assumptions are reported as sensitivities, not optimized.

\begin{table}[htbp]\centering\small
\caption{JKP reference replay, January 2000--December 2025. Annualized statistics are based on monthly excess-return accounting with modeled deductions. Utility gains are annualized quadratic-utility differences from always on. Switches are implemented changes, excluding a final-month decision that would require an unobserved next month.}\label{tab:jkp}
\resizebox{\linewidth}{!}{\input{tables/jkp.tex}}
\end{table}

The certified cumulative, retirement-only, and reversible rules make zero switches. They exactly reproduce always-on returns under each of the three cost specifications. The reference always-on Sharpe proxy is 0.475, with 6.59\% annualized mean and 13.88\% volatility. These are properties of the chosen anchored book, not evidence that the lifecycle algorithm generates incremental alpha. The certified rule's zero utility difference has a degenerate paired interval because the realized policies are identical, not because its latent economic value is known without uncertainty.

Rolling rules implement roughly 250 changes; CUSUM implements 61 and the HMM heuristic 174. Their utility differences are small. Paired circular-block intervals use 2,000 resamples of 12-month blocks of the synchronized realized utility differences. They do not rerun model selection, correct for searching across policies, or provide anytime-valid inference. In the reference replay all nonzero-policy intervals include zero.

\begin{figure}[htbp]\centering\includegraphics[width=.95\linewidth]{../figures/jkp_activity.pdf}
\caption{Active factor counts in the JKP reference replay. The certified rule remains at 153. A flat line is a substantive result, not a plotting omission.}\label{fig:activity}
\end{figure}
\begin{figure}[htbp]\centering\includegraphics[width=.95\linewidth]{../figures/jkp_utility.pdf}
\caption{Cumulative quadratic-utility differences from always on. These curves use historical factor-return proxies with modeled deductions. The certified rule has no incremental performance because it never changes allocation.}\label{fig:jkp}
\end{figure}

\subsection{Exploratory architecture diagnostics}
After observing the zero-switch reference result, we added three explicitly exploratory diagnostics: a 153-factor-only book; a 13-theme-only book; and a market-anchored 13-theme book. Their provenance is recorded in a separate manifest. No specification is selected because it performs well. These checks address whether the market anchor or large factor count alone explains inactivity.

The certified rules also make zero switches in these designs. The factor-only always-on book has an annualized Sharpe proxy of 0.466; the theme-only book has 0.573. Retirement plus revival exactly matches those benchmarks. Thus reducing the book to 13 themes does not by itself solve the power problem. The results do not support an empirical event study of certified revivals because no such events occur. Producing one from selected heuristic events and labeling it as the proposed method would be misleading.

\section{Interpretation and research implications}
The financial object is well-defined, the false-alarm accounting is explicit, and the software reproduces the economic distinction between a standalone return and a portfolio contribution. However, a valid reference method is not yet a useful active allocator in every relevant regime. Three barriers are visible in the executed evidence.

First, lifetime family-wise error control across a large strategy book is expensive. Lower-cost criteria, such as an error-rate objective tailored to repeated capital decisions, may permit more timely action, but they require a different theorem and cannot be advertised as the same 5\% lifetime guarantee. Second, clipping and fixed normalization discard economically relevant information. More efficient predictable betting and a carefully justified unbounded-score construction are natural methodological directions, not results established here. Third, an evidence threshold and a capital decision are different objects. A future economic decision layer may use uncertainty, transition persistence, and switching costs to choose fractional exposure while maintaining a separate certification layer.

A strong next experiment would compare such extensions at a \emph{matched statistical objective} and on held-out simulation designs. The current package already supplies the paired-book accounting, source audit, and adverse test cases needed for that comparison. No current result establishes minimax-optimal reversible decisions, a global regret bound, faithful counterfactual market impact, or superiority over the frontier of sequential inference.

\section{Conclusion}
An alpha's relevant state is its usefulness to a specified portfolio allocation, not merely the sign of its standalone expected return. Shadow observation makes retirement reversible, and conditional betting processes make continuous evidence collection auditable. The guarantees have precise limits: they control alarms before a directional null is violated, not arbitrary future state errors or economic regret.

The experiments show both the mechanism and its weakness. Portfolio-aware evidence can preserve a negative-return hedge and recover value after strong mean revival. It can also arrive too late to exploit covariance recovery. On 153 JKP factors the reference policy does not switch at all. The resulting research artifact is a reproducible starting point with proved statements and observed failures. Turning it into a deployable allocator requires addressing the cost of certification, not merely adding a more favorable title or a more complex prediction model.

\bigskip
\begin{thebibliography}{99}\small
\bibitem[Bhattacharyya and Ramdas(2026)]{bhattacharyya2026} Bhattacharyya, S. and A. Ramdas (2026). Change detection with conformal martingales: new optimal constructions, and suboptimality of existing methods. arXiv:2609.27179v1. \url{https://arxiv.org/abs/2609.27179}.
\bibitem[Bonacorsi(2026)]{bonacorsi2026} Bonacorsi, N. (2026). Certified Alpha Capacity: Statistical Arbitrage When Learning Takes Time. arXiv:2610.01115v2, revised 4 October 2026. \url{https://arxiv.org/abs/2610.01115v2}.
\bibitem[Choe and Ramdas(2024)]{choe2024} Choe, Y. J. and A. Ramdas (2024). Comparing Sequential Forecasters. \emph{Operations Research} 72(4), 1368--1387. Published online 17 October 2023. doi:10.1287/opre.2021.0792. \url{https://arxiv.org/abs/2110.00115}.
\bibitem[Howard et al.(2021)]{howard2021} Howard, S. R., A. Ramdas, J. McAuliffe, and J. Sekhon (2021). Time-uniform, nonparametric, nonasymptotic confidence sequences. \emph{Annals of Statistics} 49(2), 1055--1080. doi:10.1214/20-AOS1991.
\bibitem[Jensen et al.(2023)]{jkp2023} Jensen, T. I., B. Kelly, and L. H. Pedersen (2023). Is There a Replication Crisis in Finance? \emph{Journal of Finance} 78(5), 2465--2518. doi:10.1111/jofi.13249.
\bibitem[JKP Data Portal(2026)]{jkpdata} JKP Data Portal (2026). Official factor, market, and theme data; USA monthly capped-value-weighted files. Retrieved 6 October 2026. \url{https://jkpfactors.com/data}. Exact archived file hashes are in the accompanying source manifest.
\bibitem[Regis and Serra(2026)]{regis2026} Regis, M. and P. Serra (2026). Return-to-Baseline Testing via Empirically Calibrated e-processes. arXiv:2606.01960. \url{https://arxiv.org/abs/2606.01960}. Preprint.
\bibitem[Shin et al.(2022)]{shin2022} Shin, J., A. Ramdas, and A. Rinaldo (2022; revised 2023). E-detectors: a nonparametric framework for sequential change detection. arXiv:2203.03532v4. \url{https://arxiv.org/abs/2203.03532}.
\bibitem[Stephan(2026)]{stephan2026} Stephan, R. (2026). Sequential Tradeability Testing for Alpha Signals. SSRN working paper 6922558. \url{https://papers.ssrn.com/sol3/papers.cfm?abstract_id=6922558}. Citation identifies the related preprint; no priority ranking is inferred from author affiliation.
\bibitem[Waudby-Smith and Ramdas(2024)]{waudby2024} Waudby-Smith, I. and A. Ramdas (2024). Estimating means of bounded random variables by betting. \emph{Journal of the Royal Statistical Society, Series B} 86(1), 1--27. Preprint: arXiv:2010.09686. \url{https://arxiv.org/abs/2010.09686}.
\end{thebibliography}

\clearpage\appendix
\section{Proofs}\label{app:proofs}
All proofs in this manuscript are collected in this appendix. The main text contains statements, interpretations, and experimental results only.

\subsection{Proof of Proposition~\ref{prop:marginal}}
\begin{proof}
Expanding the quadratic utility gives
\[
 u_\gamma(B+A)-u_\gamma(B)=A-\gamma BA-\tfrac\gamma2 A^2.
\]
Take conditional expectations and substitute
$\E[BA\mid\F]=m_Bm_A+\Cov(B,A\mid\F)$ and
$\E[A^2\mid\F]=v_A+m_A^2$ to obtain \eqref{eq:conditionalvalue}.
The covariance can be negative and sufficiently large in magnitude for the right-hand side to be positive when $m_A<0$.
For a concrete feasible covariance structure, consider the Gaussian negative-hedge design with $q=0.2$, $w_0=0.8$, $\gamma=8$, $\mu_0=0.006$, $\mu_j=-0.003$, $f=0.0005$, variance $0.08^2$, and correlation $-0.75$. Then $m_A=-0.0007$, $m_B=0.0048$, $v_A=0.000256$, and $\Cov(A,B)=-0.000768$. Substitution gives $\E D=0.00444492>0$. This calculation includes the recurring deduction and does not rely on an unconstrained optimization argument.
\end{proof}

\subsection{Proof of Proposition~\ref{prop:optional}}
\begin{proof}
Let $\mathcal W_-$ and $\mathcal W_+$ be the feasible sets without and with access. By assumption $\mathcal W_-\subseteq\mathcal W_+$. Hence
$\sup_{w\in\mathcal W_+}\E u(R(w))\ge\sup_{w\in\mathcal W_-}\E u(R(w))$.
The comparison uses the same objective and does not introduce an unavoidable access cost. If such a cost exists, it must appear explicitly in the objective and the nesting comparison no longer has the stated form.
\end{proof}

\subsection{Proof of Theorem~\ref{thm:valid}}
\begin{proof}
Fix a strategy and episode, condition on its start-time information, and write $X_t$ for its directional evidence. Under its null,
\[
 \E[1+\lambda_tX_t\mid\F_{t-1}]
 =1+\lambda_t\E[X_t\mid\F_{t-1}]\le1,
\]
where predictability and $0\le\lambda_t<1$ ensure nonnegativity and admissibility. Thus each launched component, initialized to one and held constant before launch, is a nonnegative supermartingale until the null fails.

For the mixture, keep the mass $\pi_\ell p_\lambda$ in a dormant account before its predictable launch and then invest it in the corresponding component. The total initial capital is
$\sum_{\ell\ge1}\pi_\ell\sum_\lambda p_\lambda=1$ because
$1/[\ell(\ell+1)]=1/\ell-1/(\ell+1)$.
After $L_t$ launches the dormant account is $1/(L_t+1)$. Nonnegative mixtures preserve the conditional supermartingale property by monotone convergence. Expiry removes nonnegative capital and therefore adds only a downward jump. The resulting process is precisely \eqref{eq:mixture}, with initial value one.

Let $\nu$ be the first time $\E[X_t\mid\F_{t-1}]>0$. This conditional expectation is $\F_{t-1}$-measurable. Stop updating just before that first violation, and freeze the process thereafter. Equivalently, include only updates for which all conditional null restrictions through that update are satisfied. This predictable stopping convention produces a nonnegative supermartingale on the entire horizon. Before $\nu$, it agrees with the actually monitored process. The conditional maximal inequality consequently gives
\[
 \Prob\left(\exists t<\nu:E_t\ge1/\alpha_{j,k}\mid\F_{\sigma_{j,k}-1}\right)
 \le\alpha_{j,k}.
\]
If an episode never begins, its crossing event is empty. If it begins at a data-dependent time, the same calculation applies conditional on the information available at that start. A fresh test is funded by its designated budget, not by recycling the previous budget.

An accepted switch before $\nu_{j,k}$ is a subset of the corresponding crossing event. The atomic rule can suppress or defer a switch but cannot enlarge that event. Taking expectations and applying a countable union bound over strategies and episodes gives
\[
 \Prob\left(\bigcup_{j=1}^M\bigcup_{k\ge1}
       \{\text{accepted switch before }\nu_{j,k}\}\right)
 \le\sum_j\sum_k\frac{\delta}{Mk(k+1)}=\delta.
\]
No independence across these events is used. The proof ceases to apply to a ``false switch'' definition that includes a later economically reversed state after the same episode's null has already failed. That distinction is why the valid-prefix qualification is part of the theorem.
\end{proof}

\subsection{Proof of Theorem~\ref{thm:delay}}
\begin{proof}
Consider the selected component and suppress its indices. Since $\lambda\le\Delta/2\le1/2$ and $|X_t|\le1$, the inequality $\log(1+y)\ge y-y^2$ for $|y|\le1/2$ yields
\[
 \E[\log(1+\lambda X_t)\mid\F_{t-1}]
 \ge\lambda\Delta-\lambda^2
 \ge\frac{3\Delta^2}{16}.
\]
The last inequality follows by minimizing $\lambda\Delta-\lambda^2$ over $[\Delta/4,\Delta/2]$.
The range of a log increment has length
$\log[(1+\lambda)/(1-\lambda)]\le4\lambda$.
The martingale Hoeffding bound therefore implies, with probability at least $1-\beta$,
\[
 \log K_n\ge\frac{3n\Delta^2}{16}
            -2\lambda\sqrt{2n\log(1/\beta)}
 \ge\frac{3n\Delta^2}{16}
            -\Delta\sqrt{2n\log(1/\beta)}.
\]
Under the stated sample-size condition, $n\ge128\Delta^{-2}\log(1/\beta)$, so the second term is at most $n\Delta^2/8$. Consequently
$\log K_n\ge n\Delta^2/16\ge A$.
Because the component has mixture mass $\pi_\ell p_\lambda$, this implies
$\pi_\ell p_\lambda K_n\ge1/\alpha_{j,k}$, hence the full nonnegative mixture crosses the same boundary.

If a switch has already occurred, the stated conclusion is satisfied. Formally, the drift argument may be extended after an earlier switch by artificial bounded positive-drift variables; the event of no earlier switch is unaffected. Component expiry or numerical truncation before the needed evidence is accumulated would break this lower-bound argument, which is why the theorem excludes those cases. Competing accepted switches may defer the candidate's actual allocation change even after its evidence crosses; the theorem explicitly distinguishes crossing from acceptance.
\end{proof}

\subsection{Proof of Proposition~\ref{prop:lower}}
\begin{proof}
Let $A$ denote the event of a positive decision by observation $n$. A rule that stops earlier is still measurable with respect to the full $n$-observation experiment. Independence and the chain rule imply
\[
 \KL(P_\Delta^{\otimes n}\Vert P_0^{\otimes n})
 =n\KL(P_\Delta\Vert P_0).
\]
The data-processing inequality for the map sending the data to $\ind_A$ gives a lower bound of
$\kl(P_\Delta(A),P_0(A))$.
Since $P_\Delta(A)\ge1-\beta>\alpha\ge P_0(A)$, monotonicity of binary relative entropy in this region gives
$n\KL(P_\Delta\Vert P_0)\ge\kl(1-\beta,\alpha)$.
Here
\[
 d(\Delta)=\KL(P_\Delta\Vert P_0)
 =\frac{1+\Delta}{2}\log(1+\Delta)
  +\frac{1-\Delta}{2}\log(1-\Delta).
\]
Its derivative is $d'(x)=\operatorname{atanh}(x)\le x/(1-x^2)\le4x/3$ on $[0,1/2]$. Integrating from zero yields $d(\Delta)\le2\Delta^2/3\le\Delta^2$, sufficient for the second inequality.
\end{proof}

\subsection{Proof of Proposition~\ref{prop:clip}}
\begin{proof}
The pointwise identity
$|D-s\clip(D/s,-1,1)|=(|D|-s)_+$
and conditional Jensen's inequality give the first bound. For $|D|>s$,
$(|D|-s)_+\le |D|\le |D|^{1+\eta}/s^\eta$; the same upper bound also holds on the complement. Conditional expectation gives the stated remainder bound. A sign conclusion requires the retained signal to exceed this remainder, which is not imposed in the empirical experiment.
\end{proof}

\subsection{Proof of Proposition~\ref{prop:impossibility}}
\begin{proof}
Let $p$ be the probability of choosing active. Since the two laws agree on the information used by the decision, $p$ is the same under both, including when the rule uses independent randomization. Under the law with positive next-period incremental utility the error probability is $1-p$; under the law with negative incremental utility it is $p$. Their maximum is at least $1/2$. Neither past predictive success nor an optional-stopping argument changes the common-information premise.
\end{proof}


\section{Exact simulation specification}\label{app:dgp}
The core drift is 0.006, base standard deviation 0.08, risk-aversion coefficient 8, and reference candidate allocation $q=0.2/M$. For the one-candidate 1,800-observation experiment, change dates are 600 and 1,200. Values apply before the recurring deduction.
\begin{center}\small
\begin{tabular}{lcc}
\toprule Design & Candidate mean & Common-factor loading \\\midrule
Permanent death & $0.045\to-0.045$ & $0.2$ \\
Mean revival & $0.045\to-0.045\to0.045$ & $0.2$ \\
Negative-return hedge & $-0.003$ & $-0.75$ \\
Positive redundant & $0.006$ & $0.9$ \\
Correlation revival & $0.003$ & $-0.8\to0.8\to-0.8$ \\
Gradual correlation & $0.003$ & $-0.8\cos(2\pi t/T)$ \\
Repeated correlation & $0.003$ & Six alternating blocks of $\pm0.8$ \\\bottomrule
\end{tabular}
\end{center}
For multiple candidates, candidate $j$'s loading and drift arrays are shifted by $(j\bmod4)\lfloor T/12\rfloor$ positions. The resulting four phase groups produce cross-sectional dependence through the common shock and idiosyncratic shocks within each candidate. This is a designed regime system, not a fitted JKP factor structure.

Student innovations are independent $t_5\sqrt{3/5}$, giving unit variance. For GARCH stress tests the common predictable multiplier obeys
$h_{t+1}=0.05+0.08h_t\varepsilon_{0,t}^2+0.87h_t$, starting at one, with $s_t=0.08\sqrt{h_t}$. The autoregressive stress test instead uses
$F_t=0.25F_{t-1}+\sqrt{1-0.25^2}\varepsilon_{0,t}$.
The evaluator conditions on $F_{t-1}$, so predictable means and variances are correct rather than treating autocorrelation as independent noise.

The null experiments simulate scores directly. In the independent design,
\[Z_{j,t}=\tanh(0.6\varepsilon_{0,t}+0.8\varepsilon_{j,t}).\]
In the predictable-volatility design, the argument is multiplied by
$\clip(0.4+0.8|Z_{j,t-1}|,0.25,2)$.
Both have conditional mean zero by symmetry and therefore satisfy the directional nulls with the indifference band. The paths share a common shock across strategies, so the family-wise check is not an independence calculation.

\subsection{Implementation conventions}
The code uses a ring of $H/d=30$ restart slots and seven betting fractions, with vectorization across Monte Carlo paths and candidate strategies. The minimum memory is proportional to $M(H/d)|\Lambda|$. A restart overwrites expired wealth without normalization. Large component wealth is capped at $10^{100}$ as a downward numerical operation; the false-alarm proof survives, while the delay statement requires that this cap not prevent the relevant crossing.

The CUSUM heuristic updates $C_t=\max(0,C_{t-1}+X_t-0.05)$ and switches at eight. The HMM heuristic uses a 2\% symmetric transition probability, score emission means $\pm0.20$, score standard deviation $0.60$, and posterior threshold 0.95. It is a fixed diagnostic filter, not estimated from the simulated latent regime or tuned on evaluation results. Rolling rules average 60 scores and use the same $h=0.02$ directional threshold. Rolling buffers are not cleared after switches, a source of possible chattering that is retained in the stated comparator.

\subsection{Executed run ledger}
\begin{center}\small
\begin{tabular}{lrrl}
\toprule Suite & Scenario-path evaluations & Periods & Purpose \\\midrule
Main & 2,800 & 1,800 & 7 designs, 2 innovation laws, 200 paths each \\
Null & 4,000 & 720 & 2 known-null designs, 2,000 paths each \\
Dependence & 800 & 1,200 & 4 designs, GARCH and AR, 100 paths each \\
Scaling & 120 & 900 & 4 strategy counts, 30 paths each \\
Weak signal & 200 & 720 & Correlation magnitude divided by four \\
Delay & 1,500 & 1,800 & 3 designs, 500 paths each \\\bottomrule
\end{tabular}
\end{center}
The ledger sums to 9,420 scenario-path evaluations. Some suite specifications share base random seeds and common random numbers; the total must not be interpreted as 9,420 mutually independent financial worlds. Independent-replication standard errors apply within each reported setting. Each suite has a manifest with executed flags, parameters, seeds, and software information. Raw method-level records and deterministic replication-zero traces are retained.

\section{Additional executed results}\label{app:results}
\begin{table}[htbp]\centering\small
\caption{JKP reference paired-block intervals. Utility differences are annualized and multiplied by $10^4$. Circular 12-month blocks, 2,000 resamples. These are descriptive replay intervals, not anytime-valid confidence sequences.}
\input{tables/jkp_ci.tex}
\end{table}
\begin{table}[htbp]\centering\scriptsize
\caption{Dependence stress tests. Utility gain over always on and its Monte Carlo standard error, both multiplied by $10^4$; 100 paths per design.}
\input{tables/dependence.tex}
\end{table}
\begin{table}[htbp]\centering\small
\caption{Scaling diagnostic for correlation revival. Utility differences and MC standard errors are multiplied by $10^4$; 30 paths per strategy count. This is not a global-oracle regret comparison.}
\input{tables/scaling.tex}
\end{table}
\begin{table}[htbp]\centering\scriptsize
\caption{Exploratory JKP book-architecture checks added after the reference replay. Utility gain is relative to always on within each design. All results, including zero certified switches, are retained.}
\resizebox{\linewidth}{!}{\input{tables/diagnostics.tex}}
\end{table}
\clearpage
\section{Reproducibility and audit protocol}\label{app:reproduce}
The source package is organized as \texttt{src/alpha\_lifecycle} for portfolio accounting and monitors; \texttt{experiments} for Monte Carlo, delay, JKP replay, diagnostics, and reporting; \texttt{tests} for implementation checks; \texttt{results} for executed outputs; and \texttt{paper} for this manuscript. A single reproduction script executes the documented suites and rebuilds all tables and figures. The official JKP download is kept separate from the raw-data-free public repository.

The audit preserves four distinctions. Simulated returns are labeled simulated and are never described as JKP calibration. Historical JKP results use actual official return archives and have source hashes. Statistical null calibration uses data-generating processes for which the conditional-score null is known. Economic labels and delay diagnostics use the simulator's true raw conditional moments and do not inherit a bounded-score error guarantee by implication.

Implementation checks cover dormant-capital accounting, episode spending, paired costs and fixed-slot weights, covariance consistency, negative-return hedges, positive-return redundant allocations, prefix invariance, and atomic switching. A strong-signal threshold-crossing test checks that the implementation can act rather than being identically inert. Unit tests and Monte Carlo results complement but do not replace the mathematical proofs.

The current empirical evidence does not justify a publication-ready claim of profitable lifecycle timing. In particular, the code supplies no synthetic placeholder ``JKP result,'' no fabricated certified revival events, and no hindsight-selected representative simulation path. Tables and figures are generated from the stored executed results. Any later tuning should create a new manifest and preserve the current adverse results as a baseline.

\end{document}
